% Supplement subsections typed below as "Section~S<n>.<m>"; build.sh checks each number against the built supplement.
% supplement-ref supp:layer=S2.1
% supplement-ref supp:worked-example=S2.2
% supplement-ref supp:borda=S1.8
% supplement-ref tab:settings=S1
\section{Materials and Methods: The Ranking Layer}\label{sec:ranking-layer}

The ranking layer is the basic unit of ArrowFlow. It receives a ranking and measures how far that ranking lies from each ranking in a bank of learned rankings. Then it outputs the bank sorted by that distance. Learning reorders the learned rankings by collecting votes for where their items should sit, and no derivative is computed anywhere in the ranking layers. This section defines the layer, the training rule and the readout exactly.

\subsection{Notation}\label{sec:notation}

A \emph{ranking} of a vocabulary of $V$ items is a permutation $\pi$ that lists every item exactly once. We write $\pi[p]$ for the item at position $p\in\{1,\dots,V\}$ and $\operatorname{pos}_\pi(v)$ for the position of item $v$. The two maps are inverses, and $\operatorname{rev}(\pi)$ is $\pi$ read backwards. Worked examples list items by position, so $\pi=[C,A,E,B,D]$ means $\pi[1]=C$. Every item carries a distinct integer identifier (ID), and every sort in the method breaks exact ties by ascending ID (Section~S2.1). Table~\ref{tab:notation} collects the symbols used from here on.

\begin{table}[!htbp]
\caption{\textbf{Notation.} The layer superscript $(\ell)$ is dropped when a single layer is discussed.}
\label{tab:notation}
\centering
\small
\begin{tabular}{@{}l p{0.62\textwidth}@{}}
\toprule
Symbol & Meaning \\
\midrule
$\pi$, $\pi[p]$, $\operatorname{pos}_\pi(v)$, $\operatorname{rev}(\pi)$ & a ranking; its item at position $p$; the position of item $v$; the reversed ranking \\
$\pi_x$, $\pi'$ & the encoded input ranking; the output ranking of a layer \\
$V$ & vocabulary size of a layer, the number of items it ranks \\
$\ell$, $L$ & layer index and output layer; layer $\ell$ has vocabulary size $V_\ell$ and $N_\ell$ filters, and $V_{\ell+1}=N_\ell$ \\
$e$, $C$ & encoded vocabulary size, $V_1=e$; number of classes, $N_L=C$ \\
$r_j$, $N$ & ranking filter $j$ and the number of filters in a layer \\
$m_{\pi\to\tau}$, $m_j$ & motion of $\pi$ toward $\tau$, \eqref{eq:motion}; $m_j=m_{\pi_x\to r_j}$ \\
$\mathcal{M}^{(\ell)}$, $\mathcal{M}^{(\ell)}_j$ & signal list passed to layer $\ell$, sorted by $(-|\cdot|,\mathrm{ID})$; the entry of filter $j$ \\
$d_F$, $D_j$ & footrule distance, \eqref{eq:footrule}; $D_j=d_F(\pi_x,r_j)$ \\
$\pi^{(\ell)}$, $\hat y$ & output ranking of layer $\ell$, with $\pi^{(0)}=\pi_x$; class $\hat y=\pi^{(L)}[1]$ of the prototype readout \\
$A_j$, $A_j[p,k]$, $a$, $s_p$ & accumulator of filter $j$ (row $p$ is the item at position $p$, column $k$ a position); signed vote weight; mean voted position of row $p$ \\
$\eta$, $\rho$, $c$, $p_c$ & learning rate, selected fraction, signal scale, acceptance probability of a correct example \\
$w$, $g$ & sample weight (one throughout this paper); gate value in $\{0,1\}$ \\
$T$, $t$, $B$, $\nu$ & number of iterations, iteration index, batch size, validation fraction \\
$u$, $\tilde u$ & mean relayed motion and the scaled signal passed to the layer below, \eqref{eq:signal} \\
$h$, $\tilde h$, $q_p$, $m_y$, $m_w$ & for the encoder network of Section~\ref{sec:learned-encoder}: its scores, sorted in decreasing order into $\pi_x=\orderdesc(h)$; their targets; the target position of the coordinate at position $p$; the motions $m_j$ for $j=y$, the true class, and $j=w$, the nearest wrong class \\
\bottomrule
\end{tabular}
\end{table}

\subsection{Ranking Filters and Motions}\label{sec:filters}

A layer compares its input with each of its filters by asking how far the items must move. It holds $N$ \emph{ranking filters} $r_1,\dots,r_N$. A ranking filter is a permutation prototype, a learned permutation of the layer's vocabulary that the layer compares with its input.

The comparison rests on the \emph{motion} of one ranking $\pi$ toward another ranking $\tau$ of the same vocabulary,
\begin{equation}
m_{\pi\to\tau}[p]=\operatorname{pos}_\tau(\pi[p])-p,\qquad p=1,\dots,V,
\label{eq:motion}
\end{equation}
the signed number of positions that the item at position $p$ of $\pi$ must move to reach its position in $\tau$. A positive motion points toward the end of the ranking, and a negative one toward the front. For an input $\pi_x$ and a filter $r_j$ the motion is $m_j=m_{\pi_x\to r_j}$. Summing the absolute motions gives Spearman's footrule distance \citep{diaconis1977footrule},
\begin{equation}
d_F(\pi,\tau)=\sum_{p=1}^{V}\bigl|m_{\pi\to\tau}[p]\bigr|=\sum_{v}\bigl|\operatorname{pos}_\pi(v)-\operatorname{pos}_\tau(v)\bigr|.
\label{eq:footrule}
\end{equation}
The footrule is symmetric and takes even whole-number values from $0$ to $\lfloor V^2/2\rfloor$. It is the only distance the network computes, although the batch update minimizes squared position differences (Section~\ref{sec:theory-aggregation}). Figure~\ref{fig:forward-pass} shows the motions of one input toward three filters and the distances $D_j=d_F(\pi_x,r_j)$.

\begin{figure}[!htbp]
\centering
\resizebox{\textwidth}{!}{%
\begin{tikzpicture}[
    >=Stealth,
    arr/.style={->, thick},
    lbl/.style={font=\small},
    cell/.style={draw, minimum width=0.7cm, minimum height=0.7cm, font=\small, inner sep=0pt},
    hcell/.style={cell, fill=blue!10},
    fcell/.style={cell, fill=orange!10},
    mcell/.style={cell, fill=green!10},
]
% === LEFT HALF: Input + Filters ===
% Title
\node[font=\small\bfseries, anchor=south west] at (0.5, 1.8) {Input and Filters};
%
\node[lbl, anchor=east] at (0.9, 1.2) {Input $\pi_x$:};
\foreach \i/\v in {1/C, 2/A, 3/E, 4/B, 5/D} {
    \node[hcell] (in\i) at (0.5+\i*0.8, 1.2) {\v};
}
\node[lbl, below=0.04 of in1, font=\scriptsize] {$p{=}1$};
\node[lbl, below=0.04 of in5, font=\scriptsize] {$p{=}5$};
%
\node[lbl, anchor=east] at (0.9, -0.1) {$r_1$:};
\foreach \i/\v in {1/A, 2/B, 3/C, 4/D, 5/E} {
    \node[fcell] (f1\i) at (0.5+\i*0.8, -0.1) {\v};
}
\node[lbl, anchor=east] at (0.9, -1.0) {$r_2$:};
\foreach \i/\v in {1/C, 2/E, 3/A, 4/B, 5/D} {
    \node[fcell] (f2\i) at (0.5+\i*0.8, -1.0) {\v};
}
\node[lbl, anchor=east] at (0.9, -1.9) {$r_3$:};
\foreach \i/\v in {1/E, 2/D, 3/B, 4/A, 5/C} {
    \node[fcell] (f3\i) at (0.5+\i*0.8, -1.9) {\v};
}
% === Arrow between halves (above the rows, no overlap) ===
\draw[arr, gray, very thick] (5.3, 0.2) -- (6.8, 0.2);
\node[font=\footnotesize, black, above=2pt, align=center] at (6.05, 0.2) {compute\\motion};
% === RIGHT HALF: Motions ===
% Title
\node[font=\small\bfseries, anchor=south west] at (7.0, 2.1) {Motions $m_j[p] = \operatorname{pos}_{r_j}(\pi_x[p]) - p$};
%
\node[lbl, anchor=east] at (7.7, 1.2) {$m_1$:};
\foreach \i/\v/\c in {1/{+2}/red!20, 2/{-1}/blue!15, 3/{+2}/red!20, 4/{-2}/blue!20, 5/{-1}/blue!15} {
    \node[mcell, fill=\c, font=\small\bfseries] (m1\i) at (7.2+\i*0.8, 1.2) {\v};
}
\node[lbl, right=0.25 of m15] {$D_1{=}8$};
%
\node[lbl, anchor=east] at (7.7, 0.2) {$m_2$:};
\foreach \i/\v/\c in {1/{\phantom{+}0}/white, 2/{+1}/red!10, 3/{-1}/blue!10, 4/{\phantom{+}0}/white, 5/{\phantom{+}0}/white} {
    \node[mcell, fill=\c, font=\small] (m2\i) at (7.2+\i*0.8, 0.2) {\v};
}
\node[lbl, right=0.25 of m25] {$D_2{=}2$};
%
\node[lbl, anchor=east] at (7.7, -0.8) {$m_3$:};
\foreach \i/\v/\c in {1/{+4}/red!30, 2/{+2}/red!20, 3/{-2}/blue!20, 4/{-1}/blue!10, 5/{-3}/blue!25} {
    \node[mcell, fill=\c, font=\small\bfseries] (m3\i) at (7.2+\i*0.8, -0.8) {\v};
}
\node[lbl, right=0.25 of m35] {$D_3{=}12$};
% Formula note (below title, right-aligned)
\node[font=\footnotesize, anchor=east] at (11.8, 1.75) {$D_j = \sum_p |m_j[p]|$};
% === Output ranking ===
\node[draw, thick, rounded corners=3pt, fill=yellow!12, font=\small\bfseries, inner sep=4pt] (outbox) at (9.5,-1.8) {Output: $\pi' = [r_2,\; r_1,\; r_3]$ \;\small(closest first)};
\draw[arr, thick] (9.5, -1.2) -- (outbox.north);
\end{tikzpicture}%
}
\caption{\textbf{Ranking layer forward pass.} The input ranking $\pi_x = [C, A, E, B, D]$ is compared with three ranking filters $r_1, r_2, r_3$. The motion $m_j[p] = \operatorname{pos}_{r_j}(\pi_x[p]) - p$ is the signed displacement of the item at position $p$ toward its position in the filter (red: positive, toward the end of the ranking; blue: negative, toward the front). The footrule distance $D_j = \sum_p |m_j[p]|$ totals the displacements. The filters are ranked by distance, and $r_2$ (distance 2) is closest. The output ranking $\pi'$ is the input of the next layer.}
\label{fig:forward-pass}
\end{figure}


\subsection{Forward Pass}\label{sec:forward}

A layer's output is simply its filters, sorted from nearest to farthest, as in neural gas \citep{martinetz1993neural}. The output $\pi'$ lists the filter IDs sorted by $(D_j,j)$, closest first, with ties going to the lower ID. This output is a ranking of a new vocabulary, the $N$ filter IDs, and it is the next layer's input.

A network stacks $L$ ranking layers. Layer $\ell$ holds $N_\ell$ filters over $V_\ell$ items. The first layer ranks the $V_1=e$ items of the encoded input (Section~\ref{sec:encoding}), and every later layer ranks the filters of the layer below, so $V_{\ell+1}=N_\ell$. We write $\pi^{(0)}=\pi_x$ and $\pi^{(\ell)}$ for the output of layer $\ell$. The output layer $L$ holds $N_L=C$ filters, one per class, and its output $\pi^{(L)}$ is the \emph{class ranking}. The \emph{prototype readout} takes the class $\hat y=\pi^{(L)}[1]$, whose filter is nearest to the last hidden ranking $\pi^{(L-1)}$. The output layer only serves training. It drives the updates and makes none of ArrowFlow's predictions (Section~\ref{sec:prediction}, Table~\ref{tab:overview}).

% Method overview. Hand-written: it names parts and roles, not numbers.
\begin{table}[!htbp]
\caption{\textbf{What a view stores, and what each part is for.} A view is one encoder with the network trained on its rankings, and ArrowFlow trains $K$ of them (Section~\ref{sec:encoding}). A target-aware view puts coordinates fitted on the class labels before random ones, and a calibrated view standardizes its projected coordinates (Section~\ref{sec:encoder}). The output layer serves training and makes none of ArrowFlow's predictions.}
\label{tab:overview}
\centering
\footnotesize
\begin{tabular}{@{}p{0.19\textwidth} p{0.37\textwidth} p{0.34\textwidth}@{}}
\toprule
Part & What is fitted or stored & Role \\
\midrule
Encoder & imputation and standardization statistics and the projection matrix; a target-aware view also fits linear discriminant coordinates on the training labels, and a calibrated view stores the means and scales of its projected coordinates & turns a real-valued example into a ranking \\
\addlinespace
Hidden ranking layers & one permutation per ranking filter, reordered by position votes & produce the ranking that prediction reads \\
\addlinespace
Output layer & one permutation per class & gates the update, returns the signal from which every hidden layer learns, and scores the checkpoint \\
\addlinespace
Nearest-neighbor readout & the last hidden ranking of every training row with its label, and the chosen neighbor count and weighting & predicts the class of one view \\
\addlinespace
Vote across views & the predicted class of each view & gives the prediction of the ensemble \\
\bottomrule
\end{tabular}
\end{table}


\subsection{Position Votes and the Accumulator}\label{sec:accumulation}

Learning reorders a filter by counting votes for where each of its items should sit. Each filter $r_j$ owns an \emph{accumulator} $A_j\in\R^{V\times V}$, a table that collects these votes. Row $p$ belongs to the item $r_j[p]$ currently at position $p$, and column $k$ is a position. Between updates $A_j$ is the identity matrix, which gives every item one prior vote for its current position.

\begin{figure}[!htbp]
\centering
\resizebox{\textwidth}{!}{%
\begin{tikzpicture}[
    >=Stealth,
    arr/.style={->, thick},
    lbl/.style={font=\small},
    cell/.style={draw, minimum width=0.7cm, minimum height=0.7cm, font=\small, inner sep=0pt},
    tcell/.style={cell, fill=orange!25},
    rowlbl/.style={font=\footnotesize, anchor=east},
    colhdr/.style={font=\footnotesize},
    votelbl/.style={font=\footnotesize, anchor=east},
    note/.style={font=\scriptsize, text=black!60},
    ptitle/.style={font=\small\bfseries, align=center},
]
% ============================================================
% ROW 1:  (a) accumulator before the batch   +   (b) votes from one batch
% ============================================================
% --- (a) accumulator: identity in the current order r_j = [A,B,C,D] ---
\node[ptitle] at (1.4, 3.75) {(a) Accumulator $A_j$ before the batch};
\node[note] at (1.4, 3.3) {column $k$ = position};
\foreach \c in {1,2,3,4} {
    \node[colhdr] at (-0.35+\c*0.7, 2.55) {\c};
}
\foreach \r/\item in {1/A, 2/B, 3/C, 4/D} {
    \node[rowlbl] at (-0.1, 2.5-\r*0.7) {$p{=}\r$: \item};
    \foreach \c in {1,2,3,4} {
        \pgfmathtruncatemacro{\val}{\r==\c ? 1 : 0}
        \pgfmathtruncatemacro{\shade}{\r==\c ? 25 : 0}
        \node[cell, fill=blue!\shade] at (-0.35+\c*0.7, 2.5-\r*0.7) {\val};
    }
}
\node[note, anchor=north] at (1.0, -0.8) {row $p$: item at position $p$};
% --- plus sign ---
\node[font=\LARGE\bfseries] at (3.6, 0.9) {$+$};
% --- (b) votes from one batch ---
\node[ptitle] at (9.4, 3.75) {(b) Votes from one batch};
\node[note, text=black] at (9.4, 3.3) {a vote of weight $a$ adds $|a|$ to the column of each item's target position};
% vote 1: positive weight
\node[votelbl] at (6.7, 2.3) {vote 1 ($a_1 = {+}1$)};
\foreach \i/\v in {1/C, 2/A, 3/D, 4/B} {
    \node[tcell] at (6.55+\i*0.7, 2.3) {\v};
}
\node[font=\scriptsize, anchor=west] at (9.85, 2.3) {$\Rightarrow$ ${+}1$ at $(A,2)$, $(B,4)$, $(C,1)$, $(D,3)$};
% vote 2: negative weight, target reversed
\node[votelbl] at (6.7, 1.3) {vote 2 ($a_2 = {-}0.5$)};
\foreach \i/\v in {1/D, 2/B, 3/C, 4/A} {
    \node[tcell, fill=orange!10] at (6.55+\i*0.7, 1.3) {\v};
}
\node[font=\scriptsize, anchor=west] at (9.85, 1.3) {$a_2 < 0$: the target is reversed};
\node[votelbl] at (6.7, 0.5) {reversed ($a_2 < 0$)};
\foreach \i/\v in {1/A, 2/C, 3/B, 4/D} {
    \node[tcell] at (6.55+\i*0.7, 0.5) {\v};
}
\node[font=\scriptsize, anchor=west] at (9.85, 0.5) {$\Rightarrow$ ${+}0.5$ at $(A,1)$, $(C,2)$, $(B,3)$, $(D,4)$};
% ============================================================
% Equals sign + down arrow
% ============================================================
\node[font=\LARGE\bfseries] at (3.6, -1.35) {$=$};
\draw[arr, very thick] (3.6, -1.7) -- (3.6, -2.3);
% ============================================================
% ROW 2:  (c) after the batch  ->  (d) row means and reorder  ->  (e) reset
% ============================================================
% --- (c) accumulator after the batch (identity prior + both votes) ---
\node[ptitle] at (1.4, -2.3) {(c) After the batch};
\foreach \c in {1,2,3,4} {
    \node[colhdr] at (-0.35+\c*0.7, -3.0) {\c};
}
\node[rowlbl] at (-0.1, -3.65) {$p{=}1$: A};
\node[cell, fill=orange!40] at (0.35, -3.65) {1.5};
\node[cell, fill=orange!25] at (1.05, -3.65) {1};
\node[cell] at (1.75, -3.65) {0};
\node[cell] at (2.45, -3.65) {0};
%
\node[rowlbl] at (-0.1, -4.35) {$p{=}2$: B};
\node[cell] at (0.35, -4.35) {0};
\node[cell, fill=orange!25] at (1.05, -4.35) {1};
\node[cell, fill=orange!12] at (1.75, -4.35) {0.5};
\node[cell, fill=orange!25] at (2.45, -4.35) {1};
%
\node[rowlbl] at (-0.1, -5.05) {$p{=}3$: C};
\node[cell, fill=orange!25] at (0.35, -5.05) {1};
\node[cell, fill=orange!12] at (1.05, -5.05) {0.5};
\node[cell, fill=orange!25] at (1.75, -5.05) {1};
\node[cell] at (2.45, -5.05) {0};
%
\node[rowlbl] at (-0.1, -5.75) {$p{=}4$: D};
\node[cell] at (0.35, -5.75) {0};
\node[cell] at (1.05, -5.75) {0};
\node[cell, fill=orange!25] at (1.75, -5.75) {1};
\node[cell, fill=orange!40] at (2.45, -5.75) {1.5};
% --- arrow (c) -> (d) with the row-mean formula ---
\draw[arr, very thick] (3.25, -4.7) -- node[above=1pt, font=\scriptsize] {$s_p = \dfrac{\sum_k k\,A[p,k]}{\sum_k A[p,k]}$} (5.25, -4.7);
% --- (d) row means and reorder ---
\node[ptitle] at (7.1, -2.2) {(d) Row means\\and reorder};
\node[lbl, anchor=west] at (5.45, -3.55) {$s_1 = 1.4$ \quad (A)};
\node[lbl, anchor=west] at (5.45, -4.1) {$s_2 = 3.0$ \quad (B)};
\node[lbl, anchor=west] at (5.45, -4.65) {$s_3 = 2.0$ \quad (C)};
\node[lbl, anchor=west] at (5.45, -5.2) {$s_4 = 3.6$ \quad (D)};
\node[font=\scriptsize, anchor=west] at (5.45, -5.75) {sort by $s_p$, ties by item ID};
\node[draw, rounded corners=3pt, fill=green!12, font=\small, inner sep=5pt] at (7.1, -6.4) {$r_j \leftarrow [A, C, B, D]$};
% --- arrow (d) -> (e) ---
\draw[arr, very thick] (8.7, -4.7) -- node[above=1pt, font=\scriptsize] {reset} (9.4, -4.7);
% --- (e) reset to identity in the new order ---
\node[ptitle] at (12.45, -2.2) {(e) Reset to identity\\in the new order};
\foreach \c in {1,2,3,4} {
    \node[colhdr] at (11.05+\c*0.7, -3.0) {\c};
}
\foreach \r/\item in {1/A, 2/C, 3/B, 4/D} {
    \node[rowlbl] at (10.9, -2.95-\r*0.7) {$p{=}\r$: \item};
    \foreach \c in {1,2,3,4} {
        \pgfmathtruncatemacro{\val}{\r==\c ? 1 : 0}
        \pgfmathtruncatemacro{\shade}{\r==\c ? 25 : 0}
        \node[cell, fill=blue!\shade] at (11.05+\c*0.7, -2.95-\r*0.7) {\val};
    }
}
\node[note, anchor=north] at (12.45, -6.15) {next batch: $r_j = [A,C,B,D]$};
\end{tikzpicture}%
}
\caption{\textbf{Accumulation and reorder of one ranking filter.} (a)~The accumulator $A_j$ of the filter $r_j = [A,B,C,D]$ starts as the identity in the current order: row $p$ belongs to the item at position $p$ and column $k$ is a position. (b)~Two votes from one batch. A vote of weight $a$ toward a target ranking adds $|a|$ to the column of each item's position in the target~\eqref{eq:vote}; a negative weight reverses the target first. (c)~The accumulator after the batch. (d)~The mean voted position $s_p$ of every row~\eqref{eq:rowmean}. Sorting the items by $s_p$, with ties broken by item ID, gives the new filter $[A,C,B,D]$. (e)~The accumulator is reset to the identity in the new order before the next batch.}
\label{fig:accumulation}
\end{figure}


A training example casts position votes with a signed weight $a$ toward a target ranking $\tau$. A positive vote pulls the filter toward $\tau$. A negative vote pulls it toward the reversal of $\tau$, which acts as a push away. The target actually used is therefore $\tilde\tau=\tau$ when $a>0$ and $\tilde\tau=\operatorname{rev}(\tau)$ when $a<0$, and every item of the filter votes for its position in $\tilde\tau$:
\begin{equation}
A_j\bigl[p,\operatorname{pos}_{\tilde\tau}(r_j[p])\bigr]\leftarrow A_j\bigl[p,\operatorname{pos}_{\tilde\tau}(r_j[p])\bigr]+|a|,\qquad p=1,\dots,V.
\label{eq:vote}
\end{equation}
A zero weight adds nothing. The vote also returns $\operatorname{sign}(a)\,m_{r_j\to\tau}$, the motion of the filter toward the unreversed target, negated when the vote pushes away. The training procedure passes this motion to the layer below.

After all the votes of a batch are in, each row's mean voted position
\begin{equation}
s_p=\frac{\sum_{k=1}^{V}k\,A_j[p,k]}{\sum_{k=1}^{V}A_j[p,k]},
\label{eq:rowmean}
\end{equation}
reorders the filter. The items $r_j[p]$ are sorted by $(s_p,\mathrm{ID})$, with the row means compared as computed (Table~S1). A filter that receives no votes keeps its order. Then $A_j$ is reset to the identity in the new order. So the accumulator keeps no evidence from one batch to the next, and only the updated filter orders carry earlier training forward. Figure~\ref{fig:accumulation} traces one update. Ranking items by their mean voted position is a Borda count, a voting rule that scores each item by the positions it receives. Section~\ref{sec:arrow} develops this reading.

% No \FloatBarrier here: while Figure 3 was still waiting for its float page, the barrier forced a page break after the
% last three lines of this subsection and left a nearly empty page. Without it, Figure 3 still prints between the start
% of Section 3.4 and the start of Section 3.5; recheck that placement after any edit that moves text in this section.
\subsection{Supervised Training}\label{sec:training}

Training turns classification errors into votes, layer by layer from the top. Each iteration draws a batch and runs the forward pass with the filters held fixed. Then it updates the layers from the output layer down. In outline, every misclassified example votes to move the filter of its class toward the ranking that the example presents to the output layer, and so does a small random share of the correctly classified ones. The items of the output layer are the filters of the layer below, so the motion returned by this vote has one entry per hidden filter. A positive entry means that the class filter ranks that hidden filter earlier than the example's ranking does. That hidden filter should come nearer to its input, so it votes toward the input. A filter with a negative entry votes toward the reversed input, which pushes it away. Algorithm~\ref{alg:training} gives the complete procedure for one network. Section~S2.2 traces it by hand on a four-item example with a tied distance and a gated batch.

\begin{algorithm}[!htbp]
\scriptsize
\caption{Training one ArrowFlow network and fitting its readout}
\label{alg:training}
\begin{algorithmic}[1]
\Require rankings $\mathcal{D}=\{(\pi,y,w)\}$ of $V_1$ items with class labels and sample weights; widths $N_1,\dots,N_{L-1}$ and $N_L=C$; iterations $T$; batch size $B$; initial rate $\eta$; fraction $\rho$; scale $c$; probability $p_c$; validation fraction $\nu$; schedule (multistep or geometric); output-layer flag (trainable or frozen); readout candidates (neighbor counts and weightings); augmentation switch; seed; the data are nonempty, $L\ge2$, $C\ge2$, $V_1\ge2$, $T\ge0$, $\eta>0$, $c\ge0$, $w\ge0$, $0\le\rho\le1$, $0\le p_c\le1$, $0\le\nu<1$ with $\mathcal{D}_{\mathrm{tr}}$ nonempty, and $T$ divisible by four under the multistep schedule
\State seed the random generator; initialize every filter $r^{(\ell)}_j$ as a uniform random permutation of its vocabulary and every accumulator $A^{(\ell)}_j\gets I$
\State $\mathcal{D}_{\mathrm{tr}}\gets\mathcal{D}$; $\mathcal{D}_{\mathrm{val}}\gets\emptyset$; \textbf{if} $\nu>0$ \textbf{then} shuffle $\mathcal{D}$; move its first $\lfloor\nu|\mathcal{D}|\rfloor$ rankings to $\mathcal{D}_{\mathrm{val}}$ and keep the rest as $\mathcal{D}_{\mathrm{tr}}$
\State \textbf{if} augmentation is on \textbf{then} extend $\mathcal{D}_{\mathrm{tr}}$ with its perturbed copies (Section~\ref{sec:augmentation})
\State $\Theta^\star\gets$ current filters; \textbf{if} $\mathcal{D}_{\mathrm{val}}$ is nonempty \textbf{then} $\epsilon^\star\gets$ prototype-readout error of $\Theta^\star$ on $\mathcal{D}_{\mathrm{val}}$ \Comment{initial checkpoint}
\For{$t=1,\dots,T$}
  \State draw $\min(B,|\mathcal{D}_{\mathrm{tr}}|)$ rankings from $\mathcal{D}_{\mathrm{tr}}$ without replacement, all of them if $B$ is unset \Comment{batch}
  \ForAll{$(\pi,y,w)$ in the batch} \Comment{forward pass with fixed filters}
    \State $\pi^{(0)}\gets\pi$
    \For{$\ell=1,\dots,L$}
      \State save $\pi^{(\ell-1)}$ as the input of layer $\ell$; $D_j\gets d_F(\pi^{(\ell-1)},r^{(\ell)}_j)$ for $j=1,\dots,N_\ell$
      \State $\pi^{(\ell)}\gets$ filter IDs sorted by $(D_j,j)$
    \EndFor
    \State $\hat y\gets\pi^{(L)}[1]$ \Comment{prototype readout}
    \State $g\gets1$ if $\hat y\ne y$; otherwise draw $\xi\sim U(0,1)$ and set $g\gets1$ if $\xi<p_c$, $g\gets0$ if not \Comment{gate}
    \State vote $a\gets g\,w\,\eta/C$ on $r^{(L)}_y$ toward $\pi^{(L-1)}$ by~\eqref{eq:vote}; $\mathcal{M}^{(L-1)}\gets$ its returned motion (zero if $a=0$)
    \State save this example's $\mathcal{M}^{(L-1)}$, keyed by the filters of layer $L-1$, and sort it by $(-|m|,\mathrm{ID})$ \Comment{output vote}
  \EndFor
  \State \textbf{if} the output layer is trainable \textbf{then} reorder every $r^{(L)}_j$ by~\eqref{eq:rowmean}; in all cases $A^{(L)}_j\gets I$ for all $j$
  \For{$\ell=L-1,\dots,1$} \Comment{hidden layers, top down}
    \ForAll{examples of the batch, with list $\mathcal{M}^{(\ell)}$ and saved input $\pi^{(\ell-1)}$}
      \State $S\gets$ the first $\lceil\rho N_\ell\rceil$ entries of $\mathcal{M}^{(\ell)}$ \Comment{selection}
      \ForAll{$j\in S$} \Comment{signed votes; the earlier relay set $d_j\gets m_{r^{(\ell)}_j\to\operatorname{rev}(\pi^{(\ell-1)})}$ if $a_j<0$}
        \State $a_j\gets2\eta\,\mathcal{M}^{(\ell)}_j/N_\ell$; vote $a_j$ on $r^{(\ell)}_j$ toward $\pi^{(\ell-1)}$ by~\eqref{eq:vote}; $d_j\gets\operatorname{sign}(a_j)\,m_{r^{(\ell)}_j\to\pi^{(\ell-1)}}$
      \EndFor
      \If{$\ell>1$}
        \State $u\gets$ mean of $\{d_j : j\in S,\ a_j\ne0\}$ aligned by item ID, zero if the set is empty
        \State $\tilde u\gets c\,V_\ell\,u/(\max_v|u_v|+10^{-5})$ \Comment{mean signal}
        \State $\mathcal{M}^{(\ell-1)}\gets$ the entries of $\tilde u$, keyed by the filters of layer $\ell-1$ and sorted by $(-|\tilde u|,\mathrm{ID})$
      \EndIf
    \EndFor
    \State reorder every $r^{(\ell)}_j$ by~\eqref{eq:rowmean}; $A^{(\ell)}_j\gets I$ for all $j$ \Comment{reorder and reset}
  \EndFor
  \State \label{alg:line-schedule}multistep: $\eta\gets0.7\,\eta$ if $t>2$ and $t\in\{T/4+1,\,T/2+1,\,3T/4+1\}$, with $T$ divisible by four; geometric: $\eta\gets0.993\,\eta$ \Comment{schedule}
  \If{$\mathcal{D}_{\mathrm{val}}$ is nonempty} \Comment{checkpoint}
    \State $\epsilon\gets$ prototype-readout error of the current filters on $\mathcal{D}_{\mathrm{val}}$; \textbf{if} $\epsilon<\epsilon^\star$ \textbf{then} $\epsilon^\star\gets\epsilon$ and $\Theta^\star\gets$ current filters
  \EndIf
\EndFor
\State $\Theta\gets\Theta^\star$ if $\mathcal{D}_{\mathrm{val}}$ is nonempty, else the current filters
\State $\mathcal{H}\gets$ the last hidden ranking under $\Theta$, with its label, of every ranking in $\mathcal{D}$ \Comment{stored hidden rankings}
\State choose the neighbor count and weighting by cross-validation on $\mathcal{H}$ \Comment{readout selection}
\State \Return $\Theta$, $\mathcal{H}$ and the chosen readout settings
\end{algorithmic}
\end{algorithm}

\paragraph{Gate.} Every example that the prototype readout misclassifies is accepted for the update. A correctly classified example is accepted only with a small probability $p_c$, and otherwise it sends no signal. The gate acts only in training.

\paragraph{Output layer.} For an accepted example of class $y$ and sample weight $w$, only the class-$y$ filter receives a vote. The vote has weight $w\eta/C$ and points toward the saved input $\pi^{(L-1)}$, and no filter of another class is pushed away. Here the learning rate $\eta$ scales the votes. Through the vote weights it also scales an equivalent gradient step (Section~S1.8). The returned motion is keyed by the hidden filters that the class filter ranks. It becomes the signal list $\mathcal{M}^{(L-1)}$ for layer $L-1$, with its entries sorted from the largest magnitude down. This motion is exactly half the gradient of the squared position distance between $\pi^{(L-1)}$ and the class filter, taken with respect to the positions of the hidden filters in $\pi^{(L-1)}$ (Section~S1.8). Nothing is differentiated to obtain it.

\paragraph{Hidden layers.} At layer $\ell$ with $N_\ell$ filters, each example selects the first $\lceil\rho N_\ell\rceil$ filters of its signal list, the ones whose signals are largest in magnitude. Here $\rho$ is the selected fraction. A selected filter $j$ whose entry in the signal list is $\mathcal{M}^{(\ell)}_j$ receives the signed vote
\begin{equation}
a_j=\frac{2\eta\,\mathcal{M}^{(\ell)}_j}{N_\ell}
\label{eq:hidden-vote}
\end{equation}
toward the example's saved input $\pi^{(\ell-1)}$ by~\eqref{eq:vote}. If a hidden layer lies below, the layer relays the motions that these votes return, $\operatorname{sign}(a_j)\,m_{r_j\to\pi^{(\ell-1)}}$. An attracting vote thus asks the layer below to bring its output closer to filter~$j$, and a repelling vote asks it to move its output away. The relayed motions are averaged, item by item, into $u$, and the signal passed down is
\begin{equation}
\tilde u=\frac{c\,V_\ell\,u}{\max_v|u_v|+10^{-5}},
\label{eq:signal}
\end{equation}
with $c$ the signal scale. All votes of a layer are cast before any of its filters is reordered by~\eqref{eq:rowmean}. Only some variants of the depth study in Section~\ref{sec:controlled} use the relay above. Every other experiment used an earlier version of it, one that did not carry the sign of the vote. After a repelling vote, the earlier relay sent a motion that asked the layer below, to first order, neither to approach the filter nor to move away from it (Section~S1.8). The difference affects only networks with two hidden layers (70 of ArrowFlow's 255 outer folds), because the output layer's votes only attract.

\paragraph{Schedule and checkpoint.} The learning rate decays on a fixed schedule (Algorithm~\ref{alg:training}, line~\ref{alg:line-schedule}). A validation fraction $\nu$ of the training rankings is held out before training. After every iteration, the validation error of the prototype readout is measured, and the filters with the lowest such error are the ones returned. Each batch update of a filter exactly minimizes a squared position objective, whose targets change with the network's predictions (Section~S1.8). When the signed vote weights of a batch sum to more than $-1$, the update also equals a gradient step followed by a move to the nearest ranking. The step descends a squared position loss of the votes, starting from the filter's current order. On position vectors this is a step of learning vector quantization, and no derivative is computed to take it (Section~S1.8). But no loss of the whole network is minimized. Section~\ref{sec:experiments} measures what the updates add over untrained filters.

\subsection{Prediction}\label{sec:prediction}

The output layer holds one filter per class, which is a single point in the space of permutations. A class may occupy several regions of that space, so ArrowFlow predicts differently. A nearest-neighbor (kNN) readout compares an example's last hidden ranking with those of all training examples. The output layer then serves only the gate, the returned signal and the checkpoint. Training thus improves one classifier, the nearest class filter, while we evaluate another, the nearest training examples. Proxy-based metric learning makes the same split \citep{movshovitz2017nofuss}.

After training, the network passes every training ranking through its trained hidden layers. It stores each resulting last hidden ranking $\pi^{(L-1)}$ with its class label. For a new ranking, the readout finds the stored rankings nearest to its last hidden ranking by the footrule, and each of these neighbors votes for its class. A vote has weight one or, under distance weighting, the inverse of the neighbor's distance. The class with the most weight wins. The neighbor count and the weighting are chosen by cross-validation over the stored training rankings, with the trained layers held fixed (Section~S2.1).
